\documentclass[12pt,a4paper]{article}
\usepackage[utf8]{inputenc}
\usepackage{amsmath,amssymb}
\usepackage{graphicx}
\usepackage[left=1.5cm,right=1cm,top=1.3cm,bottom=1.3cm]{geometry}
\usepackage{lastpage}
% Font Times New Roman (cần XeLaTeX + fontspec — uncomment nếu VPS hỗ trợ):
% \usepackage{fontspec}
% \setmainfont{Times New Roman}

\begin{document}

\begin{center}
\begin{tabular}{|p{0.4\linewidth}|p{0.35\linewidth}|p{0.25\linewidth}|}
\hline
\multicolumn{1}{|c|}{SỞ GDĐT SƠN LA} & \multicolumn{2}{c|}{ĐỀ THI CUỐI KÌ 1 NĂM 2026-2027} \\
\multicolumn{1}{|c|}{TRƯỜNG THPT CÂY MƠ} & \multicolumn{2}{c|}{Môn: TOÁN 12} \\
\multicolumn{1}{|c|}{ĐỀ CHÍNH THỨC} & \multicolumn{2}{c|}{Thời gian: 90 phút, không kể thời gian phát đề} \\
\multicolumn{1}{|c|}{\textit{(Đề thi gồm có .... trang)}} & \multicolumn{2}{c|}{} \\
\hline
\multicolumn{2}{|p{0.75\linewidth}|}{Họ và tên: ........................................................................................ \newline Số báo danh: ....................................................................................} & \multicolumn{1}{c|}{\textbf{Mã đề: 4397}} \\
\hline
\end{tabular}
\end{center}

\bigskip

\noindent\textbf{PHẦN I. Câu trắc nghiệm nhiều phương án lựa chọn.}
Thí sinh trả lời từ câu 1 đến câu 2. Mỗi câu thí sinh chỉ chọn một phương án.

\medskip

\noindent\textcolor{blue}{\textbf{Câu 1.}} Cho hàm số $y=\dfrac{x^{3}}{3} - \dfrac{3 x^{2}}{2} + 9$. Hàm số đồng biến trên khoảng nào trong các khoảng sau đây?

\noindent\begin{tabular}{p{0.25\linewidth} p{0.25\linewidth} p{0.25\linewidth} p{0.25\linewidth}}
\textbf{A.}~$(0;3)$ & \textbf{B.}~$(0;+\infty)$ & \textbf{C.}~$(-\infty;3)$ & \textbf{D.}~$(3;+\infty)$ \\
\end{tabular}

\noindent\textcolor{blue}{\textbf{Câu 2.}} Cho hàm số $y=f(x)=\dfrac{1}{2} x^2+x-6\ln x+2025$ nghịch biến trên khoảng nào sau đây?

\noindent\begin{tabular}{p{0.25\linewidth} p{0.25\linewidth} p{0.25\linewidth} p{0.25\linewidth}}
\textbf{A.}~$(-3;2)$ & \textbf{B.}~$(2;+\infty)$ & \textbf{C.}~$(0;3)$ & \textbf{D.}~$(0;1)$ \\
\end{tabular}

\noindent\textbf{PHẦN II. Câu trắc nghiệm đúng sai.}
Thí sinh trả lời từ câu 3 đến câu 4. Trong mỗi ý, thí sinh chọn đúng (Đ) hoặc sai (S).

\medskip

\noindent\textcolor{blue}{\textbf{Câu 1.}} Cho hàm số $y=\dfrac{x^2+x-5}{x-2}$.

\noindent a)~Đạo hàm $y' = \dfrac{x^2-4x+3}{(x-2)^2}$

\noindent b)~Hàm số đã cho nghịch biến trên khoảng $(1;3)$

\noindent c)~Hàm số đạt cực đại tại $x=3$

\noindent d)~Tiệm cận xiên của đồ thị hàm số đã cho là $y=x+3$

\noindent\textcolor{blue}{\textbf{Câu 2.}} 
\begin{center}
\includegraphics[width=0.6\textwidth]{images/tikz_936e340f067c.svg}
\end{center}
Cho hàm số hữu tỉ $y=f(x)=\dfrac{ax^2+bx+c}{x+a}$ có đồ thị như hình vẽ bên. Đồ thị hàm số có hai điểm cực trị là $(-1;1)$ và $(-3;-7)$.

\noindent a)~Hàm số đồng biến trên khoảng $(-2;0)$

\noindent b)~Hàm số $y=f(x)$ nghịch biến trên khoảng $(-3;-2)$

\noindent c)~Đồ thị hàm số có tâm đối xứng là $I(-3;2)$

\noindent d)~$a+b+c=10$

\noindent\textbf{PHẦN III. Câu trả lời ngắn.}
Thí sinh trả lời từ câu 5 đến câu 6.

\medskip

\noindent\textcolor{blue}{\textbf{Câu 1.}} Hàm số $y=\log_2(x^2+3x)$ nghịch biến trên khoảng $(-\infty; a)$. Giá trị lớn nhất của $a$ là bao nhiêu?
	\par

\noindent \textit{Đáp số:} $-3$

\noindent\textcolor{blue}{\textbf{Câu 2.}} Khoảng đồng biến của hàm số $y=-x^3+12 x+1$ có nhiều nhất bao nhiêu số nguyên?
	\par

\noindent \textit{Đáp số:} $3$

\noindent\textbf{PHẦN IV. Câu tự luận.}
Thí sinh trả lời từ câu 7.

\medskip

\noindent\textcolor{blue}{\textbf{Câu 1.}} Xét tính đơn điệu của hàm số $y=\dfrac{x^2+x+4}{x+1}$.

\noindent\textcolor{blue}{\textbf{Câu 2.}} Tìm các khoảng đồng biến, nghịch biến của hàm số $y=2x^3-3x^2+3$.

\begin{center}
\textbf{------------ HẾT ------------}
\end{center}
\label{lastpage}


\newpage
\begin{center}
{\Large\textbf{ĐÁP ÁN VÀ LỜI GIẢI}}
\end{center}

\noindent\textcolor{blue}{\textbf{Câu 1.}}

\noindent \textit{Đáp án: \textbf{D}}
\noindent\textbf{Lời giải. } Ta có	$y'=x^{2} - 3x$, 
	$y'=0 \Leftrightarrow x^{2} - 3x = 0\Leftrightarrow \left[\begin{array}{l}x=0\\x=3.\end{array}\right.$\\

	Bảng xét dấu $y'$.
	
\begin{center}
\includegraphics[width=0.6\textwidth]{images/tikz_f98b104c7b05.svg}
\end{center}

	Từ bảng xét dấu suy ra
	
\begin{itemize}
  \item Hàm số đồng biến trên các khoảng $(-\infty;0)$ và $(3;+\infty)$.
  \item Hàm số nghịch biến trên khoảng $(0;3)$.
\end{itemize}


\noindent\begin{tabular}{p{0.25\linewidth} p{0.25\linewidth} p{0.25\linewidth} p{0.25\linewidth}}
\textbf{A.}~$(0;3)$ & \textbf{B.}~$(0;+\infty)$ & \textbf{C.}~$(-\infty;3)$ & \textbf{D.}~$(3;+\infty)$ \\
\end{tabular}

\medskip
\noindent\textcolor{blue}{\textbf{Câu 2.}}

\noindent \textit{Đáp án: \textbf{D}}
\noindent\textbf{Lời giải. } Tập xác định $\mathscr{D}=(0;+\infty)$.\\

		Ta có $f'(x)=x+1-\dfrac{6}{x}=\dfrac{x^2+x-6}{x}$.\\

		Khi đó $f'(x)=0\Rightarrow x^2+x-6=0\Rightarrow \left[\begin{array}{l}x=-3\\x=2.\end{array}\right.$\\

		Bảng biến thiên
		
\begin{center}
\includegraphics[width=0.6\textwidth]{images/tikz_03d9ea39e8f2.svg}
\end{center}

		Hàm số nghịch biến trên $(0;2)$. Do đó hàm số nghịch biến trên $(0;1)$.

\noindent\begin{tabular}{p{0.25\linewidth} p{0.25\linewidth} p{0.25\linewidth} p{0.25\linewidth}}
\textbf{A.}~$(-3;2)$ & \textbf{B.}~$(2;+\infty)$ & \textbf{C.}~$(0;3)$ & \textbf{D.}~$(0;1)$ \\
\end{tabular}

\medskip
\noindent\textcolor{blue}{\textbf{Câu 1.}}

\noindent \textit{Đáp án: ĐSSĐ}
\noindent\textbf{Lời giải. } Tập xác định $\mathscr{D}=\mathbb{R} \setminus \{2\}$.\\

\begin{itemchoice}
\itemch Ta có $y' = \dfrac{(2x+1)(x-2) - (x^2+x-5)}{(x-2)^2} = \dfrac{x^2-4x+3}{(x-2)^2}$. 
\itemch Cho $y'=0 \Leftrightarrow x^2-4x+3=0 \Leftrightarrow \left[\begin{array}{l}x=1\\x=3.\end{array}\right.$ Ta có bảng biến thiên:

\begin{center}
\includegraphics[width=0.6\textwidth]{images/tikz_a71d79724c99.svg}
\end{center}

\itemch Từ bảng biến thiên, ta thấy tại $x=3$ hàm số đạt cực tiểu.
\itemch     Ta có: 
    
\begin{align*}
        &a=\lim\limits_{x\to +\infty}\dfrac{x^2+x-5}{x(x-2)}=1;\\
        &b=\lim\limits_{x\to +\infty}\dfrac{x^2+x-5}{x-2}-x=3.
    \end{align*}

    Ta cũng có $\lim\limits_{x\to -\infty}\dfrac{f(x)}{x}=1$; $\lim\limits_{x\to -\infty}[f(x)-ax]=3$.\\

    Do đó, đồ thị hàm số có tiệm cận xiên là đường thẳng $y=x+3$.
\end{itemchoice}

\medskip
\noindent\textcolor{blue}{\textbf{Câu 2.}}

\noindent \textit{Đáp án: SĐSS}
\noindent\textbf{Lời giải. } \begin{itemchoice}
			\itemch
			Dựa vào đồ thị ta thấy trên khoảng $(-2;-1)$ hàm số nghịch biến nên không thể đồng biến trên khoảng $(-2;0)$.
			\itemch
			Dựa vào đồ thị trên khoảng $(-3;-2)$ đồ thị hàm số số có hướng đi xuống từ trái qua phải, nên hàm số nghịch biến trên khoảng này.
			\itemch	
			Đồ thị có tâm đối xứng là $I(-2;-3)$.
			\itemch
			Vì đồ thị hàm số có tiệm cận đứng $x=-2$ nên $a=2$.\\

			Khi đó $y=f(x)=\dfrac{2x^2+bx+c}{x+2}$, vì đồ thị có hai điểm cực trị $(-1;1)$ và $(-3;-7)$ nên ta có hệ phương trình
			
$$\begin{cases}\dfrac{2-b+c}{-1+2}=1\\\dfrac{18-3b+c}{-3+2}=-7\end{cases}\Leftrightarrow \begin{cases}-b+c=-1\\-3b+c=-11\end{cases}\Leftrightarrow \begin{cases}b=5\\c=4.\end{cases}$$

			Vậy $a+b+c=2+5+4=11$.
		\end{itemchoice}

\medskip
\noindent\textcolor{blue}{\textbf{Câu 1.}}

\noindent \textit{Đáp số: $-3$}
\noindent\textbf{Lời giải. } Hàm số xác định khi $x^2+3x>0\Leftrightarrow \left[\begin{array}{l}x>0\\x<-3.\end{array}\right.$\\

		Tập xác định $\mathscr{D}=(-\infty;-3)\cup(0;+\infty)$.\\

		Đạo hàm $y'=\dfrac{2x+3}{(x^2+3x)\ln 2}$.\\

		Hàm số nghịch biến khi $y'<0\Rightarrow 2x+3<0\Leftrightarrow x<-\dfrac{3}{2}$.\\

		Kết hợp điều kiện xác định, hàm số nghịch biến trên khoảng $(-\infty;-3)$.\\

		Vậy giá trị lớn nhất của $a$ là $-3$.

\medskip
\noindent\textcolor{blue}{\textbf{Câu 2.}}

\noindent \textit{Đáp số: $3$}
\noindent\textbf{Lời giải. } Ta có $y'=-3x^2+12=0 \Leftrightarrow \left[\begin{array}{l}x=2\\x=-2.\end{array}\right.$\\

		Bảng biến thiên
		
\begin{center}
\includegraphics[width=0.6\textwidth]{images/tikz_89cdf7f658d7.svg}
\end{center}

		Hàm số đã cho đồng biến trên khoảng $(-2;2)$.\\

		Vậy khoảng đồng biến của hàm số trên có nhiều nhất $3$ số nguyên.

\medskip
\noindent\textcolor{blue}{\textbf{Câu 1.}}

\noindent\textbf{Lời giải. } Hàm số xác định trên $\mathbb{R}\backslash\left\{-1\right\}$.\\
$y=\dfrac{x^2+x+4}{x+1}\Rightarrow y'=\dfrac{x^2+2x-3}{\left(x+1\right)^2}$.\\
$y'=0\Rightarrow\left[\begin{array}{l} x=1\Rightarrow y=3 \\ x=-3\Rightarrow y=-5.\end{array}\right.$\\

		Bảng biến thiên:\\

		\centerline{
			
\begin{center}
\includegraphics[width=0.6\textwidth]{images/tikz_44c3639af24c.svg}
\end{center}
}
		Hàm số đồng biến trên khoảng $\left(-\infty;-3\right)$ và $\left(1;+\infty\right)$.\\

		Hàm số nghịch biến trên khoảng $\left(-3;-1\right)$ và $\left(-1;1\right)$.

\medskip
\noindent\textcolor{blue}{\textbf{Câu 2.}}

\noindent\textbf{Lời giải. } Ta có $y'=6x^2-6x=0 \Leftrightarrow \left[\begin{array}{l}x=0\\x=1.\end{array}\right.$\\

		Bảng biến thiên
		
\begin{center}
\includegraphics[width=0.6\textwidth]{images/tikz_754af7cd2ea7.svg}
\end{center}

		Vậy
		
\begin{itemize}
  \item Hàm số đã cho đồng biến trên khoảng $(-\infty;0)$ và $(1;+\infty)$.
  \item Hàm số đã cho nghịch biến trên khoảng $(0;1)$.
\end{itemize}


\medskip

\end{document}